\section{Cosmology and Hierarchy}

The mesh grows. New docks wake at its open frontier, beat after beat, and the woven region swells without bound. That growth is easy to read as nothing more than a counter, a way of marking that time has passed. This section argues for something stronger. The growth of the mesh \emph{is} a cosmic expansion, and the radial layering of the grown region \emph{is} the renormalization hierarchy of physics. Three things that look separate turn out to be one thing seen from three sides.

A \emph{dock} is the geometric cell of the mesh, a here-now, and the mesh is a weave of docks. The \emph{wake} is the growing edge where new docks are born. The \emph{beat} is the discrete tick of time. A \emph{shell} is the set of docks at a fixed graph distance from a chosen center, one ring of the weave. These four words carry the whole argument, so the reader should hold them clearly before going on.

\begin{principle}[One growth, three readings]
The same growth of the mesh is the beat (the direction of time), the spatial expansion (a de Sitter cosmos), and the renormalization tree (the hierarchy of scales). One geometry, three readings.
\end{principle}

\subsection{The expansion law}

The mesh is a hyperbolic honeycomb, the $\{3,4,3,4\}$ tessellation. Hyperbolic space has exponentially more room as you move outward, and the honeycomb inherits that. Count the docks shell by shell from a center. The number of docks in each shell does not grow like a polynomial in the radius. It grows by a roughly constant ratio from one shell to the next.

\begin{table}[ht]
\centering
\begin{tabular}{@{}ll@{}}
\toprule
shell & docks \\
\midrule
0 & 1 \\
1 & 24 \\
2 & 456 \\
3 & 8376 \\
\bottomrule
\end{tabular}
\caption{Dock count per shell, measured outward from a center of the $\{3,4,3,4\}$ mesh. Each shell is larger than the last by a roughly constant factor.}
\label{tab:shell-counts}
\end{table}

The ratio of successive shell sizes settles to $R \approx 11$. This is the signature of exponential growth, not polynomial growth. A polynomial lattice in three flat dimensions would give shells growing like $r^2$ in count, with a ratio that drifts toward $1$. The mesh instead keeps multiplying.

Read the three-dimensional spatial volume contained within radius $r$ as the cube of a scale factor $a$, in the usual cosmological way. If the volume grows by a fixed factor per shell, and one shell of growth is one beat, then the scale factor grows by a fixed exponential in the beat.

\begin{table}[ht]
\centering
\begin{tabular}{@{}ll@{}}
\toprule
quantity & value \\
\midrule
shell ratio $R$ & $\approx 11$ \\
scale factor & $a(t) \approx R^{t/3} = e^{H t}$ \\
Hubble rate $H$ & $\ln(R)/3 \approx 0.80$ per beat \\
acceleration & $a''/a = H^2 > 0$ \\
\bottomrule
\end{tabular}
\caption{The expansion read off the shell counting. The scale factor grows exponentially in the beat, so the expansion accelerates.}
\label{tab:expansion-law}
\end{table}

Because $a(t) = e^{Ht}$, the second derivative satisfies $a''/a = H^2 > 0$. The expansion accelerates. This is exactly the form of a de Sitter cosmos, a universe whose expansion is driven by a positive vacuum energy. The effective cosmological constant is fixed by the same rate.

\begin{result}[De Sitter from growth]
The mesh's eternal growth is an accelerating, de Sitter expansion with cosmological constant $\Lambda = 3 H^2$. It is not added by hand. It falls out of the mesh being hyperbolic.
\end{result}

A control sharpens this. Build the same shell count on a flat reference lattice instead of the hyperbolic mesh. The flat shells grow only polynomially, and the ratio drifts toward $1$, so no de Sitter expansion appears. That confirms the exponential expansion is a curvature effect carried by the hyperbolic geometry, not an artifact of how the counting is done.

\subsection{Three things unified}

The very same growth is read three ways, depending on which axis of the grown region you look along.

\begin{table}[ht]
\centering
\begin{tabular}{@{}ll@{}}
\toprule
reading & what it is \\
\midrule
the beat & one ring of growth is one tick of time, the direction of time \\
the expansion & the spatial volume per beat grows exponentially, a de Sitter cosmos \\
the hierarchy & the radial tree branches outward, the renormalization scale \\
\bottomrule
\end{tabular}
\caption{One growth, three readings. The beat axis gives time, the spatial shells give expansion, the radial branching gives the scale hierarchy.}
\label{tab:three-readings}
\end{table}

The direction of time and the expansion of space are not two facts that happen to line up. They are one growth, seen along the beat axis and seen across the spatial shells. The mesh does not first tick and then separately expand. The same waking of new docks at the frontier is the tick and is the expansion.

\begin{proposition}[Time and space are one growth]
The advance of the beat and the growth of spatial volume are the same event, the waking of new docks at the mesh's frontier, read along two axes. Their agreement is identity, not coincidence.
\end{proposition}

\subsection{The hierarchy of scales}

Now read the same grown region as a tree. Pick a center and look outward. Each dock at one shell has several children at the next shell out, and the number of children is the same constant ratio $R$ that governed the expansion. The radial structure is a self-similar branching with branching factor $\approx R$.

A self-similar branching tree is exactly the shape of a renormalization hierarchy. The radial depth of the tree is the logarithm of the number of docks on its boundary. Moving from the boundary inward is coarse-graining, averaging over fine detail to reach a coarser description. Moving from the center outward is refining, resolving more detail. This is the holographic renormalization group, and in tensor-network language it is a MERA, the multi-scale entanglement structure whose layers are the radial shells.

\begin{definition}[The radial tree]
The \emph{radial tree} is the grown mesh viewed from a center, with each shell a layer and each dock branching into $\approx R$ children at the next shell out. Its depth is $\approx \log(\text{boundary size})$, and its layers are renormalization scales.
\end{definition}

The radial axis is also the axis along which the couplings of physics run. A coupling at short distance, read at the twigs of the tree near the boundary, flows to a coupling at long distance, read at the trunk near the center, as you coarse-grain layer by layer. The running of the couplings is therefore not an extra dynamical law. It is radial coarse-graining along the same tree.

\begin{result}[Couplings run along the tree]
The running of the couplings emerges from radial coarse-graining, moving up the tree from twigs (short distance, the boundary) to trunk (long distance, the deep interior). The renormalization scale is the radial coordinate.
\end{result}

The separation between scales is geometric. The tree branches exponentially, so a small number of radial steps spans an exponentially large range of scales. The vast gap between the everyday scale and the deep coarse limit is the log depth of an exponentially branching tree, nothing more exotic.

\begin{table}[ht]
\centering
\begin{tabular}{@{}lll@{}}
\toprule
level & location on the tree & physics \\
\midrule
boundary & the twigs, short distance & the frontier, where physics is read \\
middle & the branches & the running couplings \\
deep interior & the trunk, long distance & the coarse, long-distance limit \\
\bottomrule
\end{tabular}
\caption{The hierarchy of scales as the spine of the radial tree.}
\label{tab:scale-hierarchy}
\end{table}

\subsection{Dilution toward peace}

A clean cosmological consequence follows from putting conservation together with eternal growth. Recall two facts about the base.

The first is conservation. A \emph{tone} is the vibe value carried by a site, one of $\{-1, 0, +1\}$, read as pain, peace, and pleasure. The total charge $Q$, the signed sum of all tones, is fixed forever. The \emph{knit}, the collide-then-stream law, conserves it exactly at every beat.

The second is growth. The dock count grows without bound, as the shell counting above showed. New docks keep waking at the frontier forever.

Put these together. The numerator $Q$ is fixed and the denominator, the dock count, grows forever. So the average charge per dock falls toward zero. The mesh dilutes globally toward the neutral tone.

\begin{principle}[Cosmic dilution]
The total charge is fixed and the dock count grows forever, so the average charge per dock falls toward zero. The cosmos dilutes globally toward PEACE, the $0$ tone. This is the heat-death tendency.
\end{principle}

This dilution is resisted only locally. A self is a bound, self-maintaining pattern that holds a pocket of the mesh far from the neutral value by pumping charge inward and shedding entropy outward. The global drift is toward peace, and living structure is the local counter-current that swims against it. The expansion is therefore not only a spreading of space. It is a slow global cooling toward the neutral value, with selves as the local resistance.

One detail of the wake keeps the conservation honest. New docks must wake at the neutral tone $0$. A newborn dock carrying nonzero tone would inject charge from nowhere and break conservation. So waking does not mint well-being. Whatever charge a region holds is redistributed from the existing weave, not created at the frontier. The freedom at the wake is freedom of relation, how the existing charge is arranged, not freedom of net charge.

\begin{proposition}[Creation is redistribution]
New docks wake at the neutral tone, so the wake cannot mint charge. All charge is redistributed from the existing mesh. The same conservation that forbids minting well-being is what makes the cosmos dilute toward peace.
\end{proposition}

\subsection{Honest caveats}

The qualitative result is robust and the quantitative one is not yet pinned. Both deserve a plain statement.

The map from the geometric growth ratio to a \emph{physical} Hubble rate requires identifying the beat and the dock with physical units of time and length. Until that identification is fixed, the rate $H \approx 0.80$ is in beat units, not seconds. The shape of the law, exponential growth meaning de Sitter expansion, is robust. The absolute number is not yet fixed to physical units.

A realistic multi-era cosmology is also still open. The real universe ran through a radiation era, then a matter era, and only later approached de Sitter. The mesh's clean de Sitter growth is the late-time behavior. The earlier eras require the back-reaction of matter on the growth, which slows it. That refinement is open work, not a contradiction with what is measured here.

The base is discrete and deterministic. The growth is a fixed branching ratio of a discrete honeycomb, not a tuned continuous expansion. The smooth de Sitter law is the emergent, coarse-grained reading of the discrete shell counting, valid once the shells are large enough to average over.

\subsection{Why it matters}

Two of the deepest structures in physics, the expansion of the cosmos and the hierarchy of scales, both fall straight out of the mesh's geometry. Nothing was added to deliver them.

The de Sitter expansion is the weave spreading at its frontier. The hierarchy is the radial tree branching from trunk to twigs. The dilution toward peace is the slow global cooling, with selves as the local resistance. Each of these is a reading of the one growing mesh.

This ties directly to the gravity of the previous section. The de Sitter horizon there carries the same thermodynamics, with temperature $T = H / 2\pi$ and cosmological constant $\Lambda = 3 H^2$, both fixed by the growth measured here. The Hubble rate of cosmology and the horizon temperature of gravity are the same number, read on the same growing mesh.

\begin{result}[One growing mesh]
The direction of time, the expanding space, and the running of the laws are one growth. Each ring of new docks is one beat. The tree never stops branching, and the same rate sets both the cosmic expansion and the horizon temperature of gravity.
\end{result}
